\chapter{植物是生产者}\label{ch:g6:plants-are-producers}

给一盆浇过水的花盆称重，种下一颗橡子，
几年后再回来称量这棵幼橡树：数千克的木材、树皮、
叶和根——而花盆里的土壤几乎只少了
一克（\cref{ex:g4:what-plants-need:tree}首次揭示了这一
意外）。一棵数千克重的树，是用什么做成的？本章
将给出答案，并点明\omterm{def:g1:plants-around-us:plant}{植物}独特的“营生”。

\section{有机物与无机物}

\begin{definition}[无机物与有机物]\label{def:g6:plants-are-producers:matter}
对生物学家而言，物质分为两类：
\begin{itemize}
  \item \emph{无机物}\index{无机物}：非生命世界中的物质
        ——水、空气中的气体、
        岩石，以及
        \cref{prop:g4:what-plants-need:needs}中溶解的\omterm{prop:g4:what-plants-need:needs}{无机盐}；
  \item \emph{有机物}\index{有机物}：
        由\omterm{def:g6:exploring-our-environment:components}{生物}制造、构成\omterm{def:g6:exploring-our-environment:components}{生物}体的物质——
        木材、叶、\omterm{def:g3:bones-and-muscles:muscle}{肌肉}、脂肪、糖、羊毛。可以燃烧，可以腐烂，
        且总能追溯到某个\omterm{def:g6:exploring-our-environment:components}{生物}制造者。
\end{itemize}
\end{definition}

\begin{example}[分拣一口袋的东西]\label{ex:g6:plants-are-producers:sorting}
散步归来，从外套口袋里取出：一颗鹅卵石——\omterm{def:g6:plants-are-producers:matter}{无机物}；一颗
榛子——\omterm{def:g6:plants-are-producers:matter}{有机物}（榛树制造了它）；一颗塑料纽扣——
出人意料，它按来源属于\omterm{def:g6:plants-are-producers:matter}{有机物}（塑料由石油制成，
石油是古代\omterm{def:g6:exploring-our-environment:components}{生物}转化后的遗骸）；一滴
雨水——\omterm{def:g6:plants-are-producers:matter}{无机物}；一只毛线手套——\omterm{def:g6:plants-are-producers:matter}{有机物}（绵羊制造了
羊毛）。
\end{example}

\section{生产者的营生}

\begin{proposition}[植物由无机物制造有机物]\label{prop:g6:plants-are-producers:produce}
绿色\omterm{def:g1:plants-around-us:plant}{植物}只以\emph{无机}原料为起点——
即土壤中的水和溶解的盐，以及取自
\emph{空气}中的一种气体——再加上\emph{光}的\omterm{prop:g3:balanced-diet:jobs}{能量}，
便能制造\emph{有机}物：自身的糖，并由糖进而形成木材、
叶和根。这种制造正是
\cref{prop:g4:what-plants-need:make}所许诺要解释的“营生”：
\omterm{def:g1:plants-around-us:plant}{植物}是\emph{生产者}\index{生产者}——唯一能
由纯粹的\omterm{def:g6:plants-are-producers:matter}{无机物}原料制造新\omterm{def:g6:plants-are-producers:matter}{有机物}的\omterm{def:g6:exploring-our-environment:components}{生物}。
叶片中的绿色物质是捕获光的
“机器”；制造过程发生在叶片中，且在白天进行。
\end{proposition}
\admitted

\begin{omfigure}
\begin{tikzpicture}[scale=0.92]
  % leaf in the middle
  \draw[thick, omMeth, fill=omMeth!20] (5.2,1.5)
    .. controls (5.6,2.4) and (6.8,2.5) .. (7.2,1.6)
    .. controls (6.8,0.8) and (5.6,0.75) .. (5.2,1.5);
  \draw[omMeth!60!black] (5.35,1.5) -- (7.05,1.58);
  \node[font=\small\bfseries, omMeth!60!black] at (6.2,2.75)
    {白天的绿叶};
  % inputs
  \node[font=\small, align=center] (water) at (1.4,0.35)
    {水和无机盐 \\（来自土壤，经由根）};
  \node[font=\small, align=center] (gas) at (1.4,1.7)
    {来自空气的一种气体};
  \node[font=\small, align=center] (light) at (1.4,3.1)
    {光能};
  \draw[->, thick, omDef] (water.east) -- (5.3,1.2);
  \draw[->, thick, omDef] (gas.east) -- (5.15,1.5);
  \draw[->, thick, omProp] (light.east) -- (5.5,2.05);
  % outputs
  \node[font=\small, align=center] (sugar) at (10.6,1.9)
    {\textbf{有机物}：糖，继而 \\ 形成木材、\\ 叶、根、种子};
  \draw[->, very thick, omMeth!60!black] (7.3,1.7) -- (sugar.west);
  % oxygen small note
  \node[font=\footnotesize, align=center] (oxy) at (10.6,0.5)
    {（同时释放氧气 \\ 进入空气——见注）};
  \draw[->, omExo] (7.25,1.2) -- (oxy.west);
\end{tikzpicture}

{\small \omterm{def:g5:food-webs:roles}{生产者}的营生：无机原料和光进入，
\omterm{def:g6:plants-are-producers:matter}{有机物}输出。只有绿色\omterm{def:g1:plants-around-us:plant}{植物}（及其绿色
亲缘类群）进行这种制造。}
\end{omfigure}

\begin{example}[证据汇总]\label{ex:g6:plants-are-producers:evidence}
前面三个观察现在可以串联起来。称量过的
花盆：橡树的千克重并非来自土壤——它们
主要来自空气和水。
\cref{ex:g4:what-plants-need:pale}中那株苍白的“囚徒”：无光则无制造——
\omterm{def:g1:plants-around-us:plant}{植物}在黑暗中，即使\omterm{def:g6:exploring-our-environment:components}{矿物质}储备充足，也会饿死。
缺\omterm{prop:g4:what-plants-need:needs}{无机盐}的薄荷（
\cref{exo:g4:what-plants-need:11}）：需要\omterm{prop:g4:what-plants-need:needs}{无机盐}，但
只需少量，并非主体。主体来自空气和水，
\omterm{prop:g3:balanced-diet:jobs}{能量}来自光，土壤只提供“调味”。
\end{example}

\begin{remark}[气体的双向交换]\label{rem:g6:plants-are-producers:gas}
叶片吸入的空气中的气体是\emph{二氧化碳}——
正是每一个\omterm{def:g4:breathing:movements}{呼吸}者呼出的那种气体
（\cref{prop:g4:breathing:exchange}）；而制造过程
向空气中释放\emph{氧气}——正是每一个
\omterm{def:g4:breathing:movements}{呼吸}者所需要的那种气体。\omterm{def:g1:plants-around-us:plant}{植物}也进行\omterm{def:g4:breathing:movements}{呼吸}，昼夜不停，与
所有\omterm{def:g6:exploring-our-environment:components}{生物}一样；但白天，正在工作的叶片的制造
远远超过其\omterm{def:g4:breathing:movements}{呼吸}。世界上可供\omterm{def:g4:breathing:movements}{呼吸}的空气，是
\omterm{def:g5:food-webs:roles}{生产者}的副产品——每一种\omterm{def:g1:animals-around-us:animal}{动物}都欠绿色\omterm{def:g1:plants-around-us:plant}{植物}的
一笔债。
\end{remark}

\section{生产者养活所有生物——并为自己储备}

\begin{proposition}[第一个环节，已阐明]\label{prop:g6:plants-are-producers:firstlink}
\Cref{prop:g3:food-chains:start}规定，每一条\omterm{def:g3:food-chains:chain}{食物链}
都始于一种\omterm{def:g1:plants-around-us:plant}{植物}；现在，这一规定有了其缘由。只有
\omterm{def:g5:food-webs:roles}{生产者}能由\omterm{def:g6:exploring-our-environment:components}{矿物质}制造\omterm{def:g6:plants-are-producers:matter}{有机物}；每一个
\omterm{def:g5:food-webs:roles}{消费者}（见\cref{def:g5:food-webs:roles}）只能
\emph{转化}其所取食的\omterm{def:g6:plants-are-producers:matter}{有机物}。因此，
每个\omterm{def:g6:exploring-our-environment:components}{生物}体内的全部\omterm{def:g6:plants-are-producers:matter}{有机物}——兔子、狐狸、秃鹰、你——
最初都是在叶片中制造出来的，而
\cref{prop:g5:food-webs:pyramid}中的食物金字塔立于绿色
底层之上，因为物质正是在此进入生命世界。
\end{proposition}
\admitted

\begin{example}[植物自己的储藏室]\label{ex:g6:plants-are-producers:pantry}
\omterm{def:g1:plants-around-us:plant}{植物}制造出超过当日所需的部分，便以
有机储备的形式储存：马铃薯\omterm{ex:g4:plants-without-seeds:bulbtuber}{块茎}与郁金香\omterm{ex:g4:plants-without-seeds:bulbtuber}{鳞茎}
（\cref{ex:g4:plants-without-seeds:bulbtuber}）、豆的
肥厚两瓣（\cref{def:g2:from-seed-to-plant:seed}）、
胡萝卜肥大的根（\cref{exo:g1:plants-around-us:10}）、
甜菜的甜味。它们每一种都是夏日的光，以
物质形式储存起来——而它们每一种，正是农民和
食用者所收获的东西。
\end{example}

\begin{method}[追溯任一食物至其叶片]\label{met:g6:plants-are-producers:trace}
取桌上的任一食物：
\begin{enumerate}
  \item 若为\omterm{def:g1:plants-around-us:plant}{植物}性物质——面包的小麦、
        马铃薯、苹果——则直接是\omterm{def:g5:food-webs:roles}{生产者}的产物：
        写出\omterm{def:g1:plants-around-us:plant}{植物}的名称及其储存器官；
  \item 若为\omterm{def:g1:animals-around-us:animal}{动物}性物质——奶酪、鸡蛋、火腿——
        沿\omterm{def:g3:food-chains:chain}{食物链}向下追溯：该\omterm{def:g1:animals-around-us:animal}{动物}吃了什么？以及
        那种食物的制造者是谁？继续追溯至\omterm{def:g1:plants-around-us:plant}{植物}；
  \item 用箭头写出\omterm{def:g3:food-chains:chain}{食物链}
        （\cref{def:g3:food-chains:chain}）；
  \item 结论始终相同：追溯路径止于
        叶片，正在光照之中。
\end{enumerate}
\end{method}

\section{练习}

\begin{exercise}[$\star$]\label{exo:g6:plants-are-producers:1}
分别给出\omterm{def:g6:exploring-our-environment:components}{矿物质}与\omterm{def:g6:plants-are-producers:matter}{有机物}的定义，并各
举两个例子。
\end{exercise}

\begin{exercise}[$\star$]\label{exo:g6:plants-are-producers:2}
写出\omterm{def:g5:food-webs:roles}{生产者}的“生产过程”：投入物、\omterm{prop:g3:balanced-diet:jobs}{能量}来源、
产物，以及这一过程发生的“车间”。
\end{exercise}

\begin{exercise}[$\star$]\label{exo:g6:plants-are-producers:3}
将下列物质分为\omterm{def:g6:exploring-our-environment:components}{矿物质}或\omterm{def:g6:plants-are-producers:matter}{有机物}：露水、蜗牛壳、蜂蜡、
花岗岩、食用油、海盐。
\end{exercise}

\begin{exercise}[$\star$]\label{exo:g6:plants-are-producers:4}
橡树增加的千克数主要来自何处，若几乎
不来自土壤？
\end{exercise}

\begin{exercise}[$\star$]\label{exo:g6:plants-are-producers:5}
工作中的叶片从空气中吸收哪种气体，并
释放出哪种气体？谁呼出前者，又需要
后者？
\end{exercise}

\begin{exercise}[$\star$]\label{exo:g6:plants-are-producers:6}
写出四种\omterm{def:g1:plants-around-us:plant}{植物}储藏器官及各器官所储存的物质，
据\cref{ex:g6:plants-are-producers:pantry}。
\end{exercise}

\begin{exercise}[$\star\star$]\label{exo:g6:plants-are-producers:7}
为何放在黑暗橱柜中的\omterm{def:g1:plants-around-us:plant}{植物}，即便\omterm{prop:g4:what-plants-need:needs}{无机盐}储备
充足，仍会饿死？哪一种投入无可替代？
\end{exercise}

\begin{exercise}[$\star\star$]\label{exo:g6:plants-are-producers:8}
为什么没有任何\omterm{def:g1:animals-around-us:animal}{动物}能成为\omterm{def:g3:food-chains:chain}{食物链}的第一个环节？
请结合\omterm{def:g5:food-webs:roles}{生产者}与\omterm{def:g5:food-webs:roles}{消费者}的区别作答。
\end{exercise}

\begin{exercise}[$\star\star$]\label{exo:g6:plants-are-producers:9}
对奶酪三明治运用\cref{met:g6:plants-are-producers:trace}，
即面包、黄油、奶酪——三条追溯路径，分别通向
三片\omterm{def:g1:plants-around-us:parts}{叶子}。
\end{exercise}

\begin{exercise}[$\star\star$]\label{exo:g6:plants-are-producers:10}
“堆肥喂养\omterm{def:g1:plants-around-us:plant}{植物}，正如干草喂养奶牛一样。”
请准确修正这句话，运用
\cref{rem:g4:what-plants-need:soil}及本章的
定义——堆肥实际提供的是什么，而
\omterm{def:g1:plants-around-us:plant}{植物}自身制造的又是什么？
\end{exercise}

\begin{exercise}[$\star\star$]\label{exo:g6:plants-are-producers:11}
密封、有光照的水草与\omterm{prop:g3:sorting-living-things:groups}{鱼类}水族箱可以维持
数月；密封的\emph{黑暗}水族箱则在数天内便告失败。请运用
\cref{rem:g6:plants-are-producers:gas}来解释
这两种结果。
\end{exercise}

\begin{exercise}[$\star\star\star$]\label{exo:g6:plants-are-producers:12}
一项古老的实验：一株柳树\omterm{prop:g2:from-seed-to-plant:cycle}{幼苗}在一只
桶中生长了五年，增重许多千克；桶中干燥的土壤只
减少了几克。实验者断言：“这棵树完全是由
水构成的。”他的结论中，哪一部分正确，
哪一部分不完整——以及他所不知道的第二种
\omterm{def:g6:exploring-our-environment:components}{矿物质}投入，现代生物学家会补充什么？
\end{exercise}

\section{问题：菜农的一年}

\begin{problem}[{周末问题——光、叶与
菜园账本}]\label{pb:g6:plants-are-producers:1}
一位菜农种植土豆、胡萝卜、豆角和
生菜，并精心记账，记录投入了什么、
产出了什么。

\medskip\noindent\textbf{第一部分——账本的谜题。}
\begin{enumerate}
  \item 菜园每年产出两吨蔬菜；
        菜农只撒施了几袋堆肥。请细想这一谜题：
        这些吨数究竟从何而来？若不是
        来自那几袋堆肥，又来自何处？（请结合\omterm{def:g5:food-webs:roles}{生产者}的作用作答。）
  \item 对产量进行分类：土豆、胡萝卜、
        豆\omterm{def:g2:from-seed-to-plant:seed}{种子}、生菜叶——分别说出每一部分
        对\omterm{def:g1:plants-around-us:plant}{植物}所起的作用（贮藏器官、工作器官、
        \omterm{def:g2:from-seed-to-plant:seed}{种子}储备）。
  \item 菜农说：“我真正的供货方是靠
        阳光供货的。”请用
        \cref{prop:g6:plants-are-producers:produce}中的
        术语加以转述。
  \item 既然堆肥并非为了吨数，菜农为何仍
        需要堆肥？（回想一下沙中的
        薄荷。）
\end{enumerate}

\medskip\noindent\textbf{第二部分——读懂季节。}
\begin{enumerate}[resume]
  \item 六月份的生菜在试验遮布下遮荫
        一周，变得发白、徒长。请说出
        这一做法重复的是哪个实验，去除了哪种投入？
  \item 土豆垄已经培土，其茂盛的
        地上部分备受称赞：“茎叶好，\omterm{ex:g4:plants-without-seeds:bulbtuber}{块茎}才好。”请解释
        这句话——地上工作的叶与
        地下膨大的贮藏器官之间有何联系？
  \item 八月的干旱使生长停滞，尽管太阳
        正烈。哪一种\omterm{prop:g4:what-plants-need:needs}{无机盐}供应失效了？为何光
        不能替代它？
  \item 九月，菜农割去并移走了
        豆株，却把根翻埋入土。邻居
        认为这是浪费。请借助\omterm{def:g6:exploring-our-environment:components}{矿物质}和\omterm{def:g6:plants-are-producers:matter}{有机物}的
        定义来为之辩护：什么被归还到
        土壤中？又是为了谁？
\end{enumerate}

\medskip\noindent\textbf{第三部分——超越菜园。}
\begin{enumerate}[resume]
  \item 菜农的猪吃掉小土豆；
        全家人吃猪肉。请写出这条\omterm{def:g3:food-chains:chain}{食物链}，并说明
        这些\omterm{def:g6:plants-are-producers:matter}{有机物}最初是在何处
        制造的？
  \item 一位顾客声称她的沙拉“主要是阳光”。
        请核查这一说法：列出沙拉叶片的
        实际成分，以及阳光的确切作用。
  \item 霜冻一夜之间杀死了最后的生菜；它们
        躺在原地腐烂，到十一月。它们的\omterm{def:g6:plants-are-producers:matter}{有机物}
        并不会消失。请预先思考
        \cref{ch:g6:decomposers-and-soil}：谁来分解它，
        并将其转化为何种状态？
  \item 用一句话概括这个菜园，句中须包含：
        \omterm{def:g5:food-webs:roles}{生产者}、\omterm{def:g6:exploring-our-environment:components}{矿物质}、\omterm{def:g6:plants-are-producers:matter}{有机物}、光。
\end{enumerate}
\end{problem}
